\honest{Don't Believe You'll Self-Deceive}{Eliezer Yudkowsky}{2009}

I do not mean to pick on Kurige, but Kurige wrote that a once blind faith had moved to ``somewhere in the vicinity of Orwellian double-think.'' If you know it is double-think, I want to ask, how can you still believe it? \nb{Kurige's next paragraph says what the word meant: holding both that God instilled a sense of right and wrong and that morality evolved, and ``These two beliefs are not contradictory.'' Orwell's doublethink is holding two contradictory beliefs, so Kurige used the word more loosely than the post reads it.}

Kurige also wrote: ``I chose to believe in the existence of God—deliberately and consciously. This decision, however, has absolutely zero effect on the actual existence of God.'' If you know your belief is not correlated with reality, I ask, how can you still believe it? \nb{No belief makes its object exist, accurate beliefs included; on the author's own account, a belief tracks reality by being caused by it. Kurige's comment compared the belief to a scientist's belief in the Higgs boson, held on evidence. The reading rests on ``chose.''}

Shouldn't the gut-level realization follow from knowing that a map has no reason to match the territory? ``Well... apparently not.'' Part of the answer ``may be'' my account in ``Moore's Paradox.'' Another part ``may just be'' that it is ``actually quite easy'' not to go from what an accurate map would say to believing it.

I realize now that when I wrote ``You cannot make yourself believe the sky is green by an act of will,'' I was not just reporting the facts. I was also trying to instill a self-fulfilling prophecy. \nb{The sentence is from ``Doublethink'' (2007), where it was offered as fact; ``Belief in Self-Deception'' repeated it four days before this post. The post's link for it points to ``Moore's Paradox.''}

So ``It may be wise'' to keep repeating that you cannot double-think, and that deep down you will know. Then, I say without a hedge, ``you will indeed be less likely to fool yourself successfully.'' I give no evidence. \nb{The advice is to adopt a belief for its effects, close to what ``Doublethink'' argued against. The difference is that this belief is meant to make itself true once held.}

``You must believe in your own inability!'' Tell yourself the effort is doomed, ``and it will be!'' Is that positive thinking or negative thinking? ``Either way, it seems like a wise precaution.''
